
\chapter{Data-to-Text Generation}

Data-to-text (D2T) generation is a branch of natural language generation in
which a system converts structured input data into natural-language text
\cite{gatt2018survey,osuji2024systematic}. Inputs may be tables, database
records, time series, or graphs. The output should communicate the source
content fluently, coherently, and without unsupported statements. D2T is thus a
mapping from a structured input representation to a linguistic output.

\section{Problem definition}
\label{sec:d2t-problem-definition}

\subsection{Data-to-Text task}

The general aim of D2T generation is to produce a textual description $y$ from
structured input $x$:
\begin{equation}
    G(x) = y,
\end{equation}
where $G$ is a data-to-text generator. 

D2T is distinguished from text-to-text generation by the nature of its source.
In machine translation, summarization, and paraphrasing, the input is already
natural language. In D2T, the generator must determine what the source data
means for the current communicative purpose, organize that content into a
discourse, and realize the resulting plan as sentences. The
required output can range from a short dialogue utterance to a multi-sentence
report or a complete document~\cite{gatt2018survey,osuji2024systematic}.

The generated text is normally expected to satisfy several requirements.
\emph{Faithfulness} concerns whether every expressed statement is supported by
the source; incorrect values and additions are faithfulness errors.
\emph{Adequacy} concerns whether the selected content is appropriate for the
communicative purpose, while \emph{coverage} concerns omissions among facts
that the task requires the text to communicate. The output should also be
grammatical, fluent, coherent, and appropriately organized.

\subsection{RDF and semantic triples}
\label{sec:rdf-semantic-triples}

The Resource Description Framework (RDF) is a graph-based data model for
representing statements about resources~\cite{w3c2014rdf11concepts}. The basic
unit of an RDF statement is a triple consisting of a subject, a predicate, and
an object. The subject identifies the entity being described, the predicate
denotes a relation or property, and the object identifies another entity or
contains a literal value. A compact representation of a statement about a
mobile phone can therefore be written as:
\begin{equation}
    (\texttt{ex:PhoneModel},\ \texttt{ex:hasBatteryCapacity},\ \texttt{"5000 mAh"}).
\end{equation}
Here, \texttt{ex:} is a namespace prefix, \texttt{ex:PhoneModel} is the
subject, \texttt{ex:hasBatteryCapacity} is the predicate, and the string
\texttt{"5000 mAh"} is the object literal. In a complete RDF representation,
resources and predicates are normally identified using globally resolvable
identifiers, while literals may carry datatypes or language tags.

A collection of RDF triples forms a directed, labeled graph: subjects and
resource objects are nodes, and predicates label the edges between them. The
same entity can participate in many relations, and an object can simultaneously
be the subject of other statements. RDF does not itself determine how the facts
should be ordered or aggregated in a text~\cite{w3c2014rdf11concepts}. It can
therefore serve as a \emph{meaning representation}: an explicit,
structured representation of content that has been selected for communication.
It preserves the selected facts
and their relations while leaving wording, sentence structure, and discourse
organization to the \emph{surface realization}: the process of converting the
meaning representation into natural language.

In this thesis, a simplified version of RDF triples are used, which can be written as
\begin{equation}
    T = \{(s_i,p_i,o_i)\}_{i=1}^{n},
\end{equation}
where $s_i$ is a subject, $p_i$ is a predicate, and $o_i$ is an object. From now on in the thesis, the name
\emph{semantic triples}, or just \emph{triples} will be used to refer to this simplified representation. 
It has the same structural role as an RDF triple but
does not require RDF namespaces, datatypes, or a complete RDF document. Here,
$T$ \emph{is} the meaning representation: there must must still be a process which decides how to order, aggregate, and phrase them. Several
valid texts can describe the same input~\cite{gatt2018survey,osuji2024systematic}.

For example, the triples \texttt{(PhoneModel, hasBatteryCapacity, 5000 mAh)}
and \texttt{(PhoneModel, hasMainCamera, 108 MP)} can be realized as one sentence
or two. Their representation alone does not prescribe either choice.

\subsection{Generation stages}

Classical descriptions of natural language generation divide D2T into a
sequence of steps~\cite{reiter1997building,gatt2018survey}:

\begin{enumerate}
    \item \textbf{Content determination and analysis} determines which source
    facts should be communicated and may derive summaries, extrema, or trends.
    For example, a time series can support the statement that a value is rising
    rather than require every observation to be listed.
    \item \textbf{Document structuring} determines the order in which selected
    facts or entities are introduced. The ordering can follow temporal,
    spatial, causal, or discourse-based relations.
    \item \textbf{Content aggregation} groups related facts into clauses,
    sentences, or paragraphs. Aggregation can reduce repetition and make the
    relations between facts explicit.
    \item \textbf{Lexicalization} selects words and phrases that express the
    predicates and values in the target language.
    \item \textbf{Referring expression generation} determines how entities
    should be mentioned, for example by using a proper name on first mention and
    a pronoun or a shorter description later in the discourse.
    \item \textbf{Surface realization} converts the resulting plan into
    grammatical natural-language sentences.
\end{enumerate}

The first two decisions determine ``what to say?'', and aggregation determines
how selected content is grouped before realization. Lexicalization, referring
expression generation, and surface realization address ``how to say it?''. The
first group can produce a meaning representation, such as an ordered set of semantic triples or a
sentence plan; the second group converts that meaning representation into text. The distinction is
useful even when all stages are implemented by a single neural model because it
identifies different sources of errors and possible intermediate artifacts
\cite{moryossef-etal-2019-step}.

\subsection{Properties and challenges}

There are many difficulties in D2T generation.
They depend on the input representation, the
number and relations of facts, the required discourse scope, and whether the
task includes content determination. A fixed schema can make a template system
practical, whereas varied or interconnected records increase the engineering or
learning burden.

Faithfulness is a particularly important challenge for learned generators. A
neural system can generate fluent text that contains an incorrect value, omits
a required fact, or adds information not supported by the source. The latter
two types of error are commonly studied as omissions and hallucinations
\cite{wiseman-etal-2017-challenges,dusek-kasner-2020-evaluating}. A rule-based
system can avoid many such errors by construction, but only when its rules
correctly cover the input and realization cases.

Another challenge is the multiplicity of valid realizations. A single set of
triples can be expressed with different word choices, sentence boundaries, and
fact orderings. A metric that rewards only agreement with one reference can
therefore penalize a valid output. For learned systems, this issue affects both
training and evaluation and motivates the use of multiple references,
source-aware metrics, and human or model-based assessment
\cite{dhingra-etal-2019-handling,novikova-etal-2017-need}.

\section{Methods}
\label{sec:d2t-methods}

\subsection{Rule-based and template-based systems}

The earliest D2T systems were built as manually engineered pipelines. A
computational linguist or system developer defined rules for selecting and
ordering facts, templates for common
sentence structures, and lexicalization procedures for inserting values into
those templates. A template such as ``$\langle$entity$\rangle$ is located in
$\langle$place$\rangle$'' can be instantiated whenever the input contains a
corresponding location relation.

The modular structure of these systems offers several advantages. The rules can
be inspected directly, unsupported facts can be excluded by construction,
the output style can be controlled precisely and the runtime generation is fast.
These properties made
templates attractive in domains such as weather forecasting, sports reporting,
and dialogue systems~\cite{reiter1997building,gatt2018survey}.

Their main limitation is the cost of domain engineering. New predicates,
entities, combinations of facts, and linguistic constructions require new rules
or templates. A finite template inventory also restricts lexical and syntactic
variation. As the domain becomes more varied, maintaining interactions between
content selection, aggregation, referring expressions, and surface realization
becomes difficult. The resulting systems can be faithful and controllable while
still sounding repetitive or unnatural.

\subsection{Statistical and neural sequence-to-sequence systems}

Statistical D2T systems learn generation decisions from aligned data rather
than defining each decision manually. Depending on the architecture, they can
model content selection, ordering, sentence planning, or surface realization.
Earlier work includes probabilistic approaches to these separate decisions,
learned alignments, and statistical template induction
\cite{van-der-lee-etal-2018-automated,gatt2018survey,osuji2024systematic}.

Compared with fully handwritten systems, statistical methods can generalize to
combinations that were not explicitly covered by a template. Their behavior is
nevertheless dependent on the quality, size, and domain coverage of the aligned
corpus. Sparse observations or poor alignments can lead to omissions, incorrect
lexicalization, or incoherent ordering.

Neural D2T systems are statistical models that formulate generation as
conditional language modeling. An encoder represents the structured input, and
a decoder generates the output one token at a time. Recurrent encoder--decoder
models with attention were among the first widely used neural architectures for
the task. At each decoding step, attention weights the encoded input records
that are most relevant to the next token
\cite{gatt2018survey,osuji2024systematic}.

The input may first be linearized into a token sequence, or it may be processed
by a specialized encoder. Graph encoders can represent RDF nodes and labeled
edges, while hierarchical table encoders can separately represent rows, fields,
and cell values. Then a decoder, such as a Copy mechanisms, allow names, dates, and numbers to be copied
from the input rather than generated from the vocabulary, which helps preserve
rare or previously unseen values~\cite{gatt2018survey,osuji2024systematic}.

Longer documents introduce additional problems. The decoder must remember which
records have already been mentioned, maintain a coherent discourse order, and
avoid repeating or omitting facts. The RotoWire data-to-document study showed
that neural models can produce fluent text while struggling with content
selection and long-range structure; in some settings, templated baselines
performed better on content-related measures~\cite{wiseman-etal-2017-challenges}.

Neural D2T systems can be organized as a modular pipeline or as an end-to-end
model. In a modular neural pipeline, separate models perform decisions such as
discourse ordering, text structuring, referring expression generation,
lexicalization, and realization. The output of one module is passed to the next
one. In an end-to-end system, the input is encoded and directly decoded into
text, with no explicitly supervised intermediate plan.

Castro Ferreira et al.~\cite{castro-ferreira-etal-2019-neural} compared neural
pipeline and end-to-end systems for RDF-to-text generation using GRU and
Transformer architectures. Their results indicate that explicit intermediate
steps can improve generation quality and generalization to unseen inputs. The
pipeline also makes individual decisions easier to analyze, although errors can
propagate between modules and the system requires intermediate annotations or
procedures for constructing them.

An intermediate design is provided by Moryossef et al.~\cite{moryossef-etal-2019-step}.
Their system first creates a symbolic text plan that specifies the division of
facts into sentences and their ordering, and then uses a neural model to
realize that plan. The symbolic planning stage improves control over coverage
and structure, while the neural realization stage supplies linguistic
flexibility.

\subsection{Transformers and pretrained models}

Transformer architectures are now widely used for D2T because self-attention
can model relationships among distant parts of the input and output. Pretrained
language models can be adapted through supervised fine-tuning, or prompt-based
generation. They often provide strong
linguistic fluency and can transfer knowledge across domains, but their learned
generation decisions are difficult to inspect
\cite{osuji2024systematic,kasner2024beyond}.

\section{Interpretable D2T approaches}
\label{sec:interpretable-d2t}

Interpretability in D2T generation concerns the ability to understand why a
particular text was produced from a particular input. An interpretable system
should make the connection between source facts, generation decisions, and
output phrases sufficiently explicit for a researcher or developer to inspect.
It is related to, but distinct from, controllability, and
auditability: an interpretable procedure may still make an incorrect decision,
but it provides an artifact through which that decision can be examined.

Recent work uses an LLM during system construction rather than at every
generation step. Warczyński, Lango, and Dušek
\cite{warczynski-etal-2024-leveraging-large} prompt an LLM to produce a
rule-based Python D2T generator. The resulting program performs runtime
generation on a CPU and can be inspected, edited, and evaluated independently
of the LLM.

Lango and Dušek~\cite{lango-dusek-2025-llm} extend this direction with a
neurosymbolic RDF-to-text framework in which multiple LLM agents collaborate to
construct and refine rule-based Python code. The agents use RDF triples and do
not require in-domain human reference texts or supervised backpropagation. The
resulting executable program remains the inspectable generation artifact and
can run without the construction model at inference time.

The method proposed in this thesis uses the same separation between program
construction and program execution, but changes the construction procedure.
Instead of following one agent interaction trajectory, it repeatedly mutates
programs, executes them on a fixed evaluator, and retains
measured alternatives. This provides basis for accepting or
rejecting program changes; its implementation is described in Chapter~4.

\section{Performance Evaluation}
\label{sec:d2t-performance-evaluation}

\subsection{Evaluation dimensions}

No single measurement captures all desirable properties of generated text.
Evaluation is therefore organized along several dimensions
\cite{gatt2018survey,novikova-etal-2017-need}. \emph{Text quality} includes
grammaticality, fluency, readability, and coherence. \emph{Faithfulness}
concerns whether the claims in the output are supported by the source data.
\emph{Coverage} concerns required source facts that were not expressed, while
\emph{non-redundancy} concerns unnecessary repetitions. Adequacy can be used
for the separate question of whether the selected content suits the
communicative purpose. Some applications additionally evaluate style,
diversity, latency, and computational cost.

Automatic procedures can be reference-based or source-aware. Reference-based
metrics compare a generated text with one or more human-written texts. They are
easy to apply when a corpus contains references, but they can penalize valid
paraphrases and cannot by themselves establish whether either a reference or a
candidate is supported by the source. Source-aware procedures compare the output
with the structured input and can target faithfulness and coverage directly.

\subsection{Reference-based automatic metrics}

\subsubsection{BLEU}

BLEU (Bilingual Evaluation Understudy) was introduced for machine translation by
Papineni et al.~\cite{papineni2002bleu}. It measures modified precision for matching
unigrams, bigrams, trigrams, and four-grams between a candidate and one or more
references. A brevity penalty reduces the score of candidates that obtain high
precision by being much shorter than the references. A simplified corpus-level
form is
\begin{equation}
    \mathrm{BLEU} = BP \cdot \exp\left(\sum_{n=1}^{N} w_n \log p_n\right),
\end{equation}
where $p_n$ is the modified precision for $n$-grams, $w_n$ are the weights, and
$BP$ is the brevity penalty.

BLEU is computationally inexpensive, deterministic, and useful for comparing systems under the
same implementation and preprocessing conditions. Its limitations are equally
important. It rewards lexical overlap rather than meaning, does not directly
measure factual correctness, and can assign a low score to a faithful
paraphrase. It can also assign a relatively high score to a text that shares
surface forms with a reference while omitting or altering source facts.

\subsubsection{METEOR}

METEOR was proposed by Banerjee and Lavie~\cite{banerjee2005meteor} to improve
the relationship between automatic and human evaluation. It aligns candidate
and reference unigrams using exact matches and, depending on the implementation,
stem and synonym matches. Precision and recall are combined in a weighted
harmonic mean, and a fragmentation penalty discourages matches that occur in a
disordered fashion. METEOR can therefore recognize some lexical variation that
BLEU cannot. It remains reference-dependent, however, and its behavior depends
on the available stemming and synonym resources.

\subsubsection{BERTScore}

BERTScore uses contextual embeddings rather than exact token identity to
compare a candidate with a reference~\cite{zhang2020bertscore}. Each token in
one text is matched to semantically similar tokens in the other text using
cosine similarity. Precision, recall, and F1 variants can be reported. Because
contextual embeddings represent related words and paraphrases more flexibly,
BERTScore can capture similarities missed by n-gram metrics. It is nevertheless
dependent on the pretrained model, tokenizer, language, and reference text, and
semantic similarity does not guarantee factual faithfulness to the source data.

\subsubsection{BLEURT}

BLEURT is a learned evaluation metric trained to predict human judgments of
generated text~\cite{sellam2020bleurt}. It combines pretrained language-model
representations with further training on synthetic perturbations and human
ratings. Unlike a direct n-gram count, BLEURT can learn that some substitutions,
omissions, or fluency problems affect quality differently. This flexibility can
improve correlation with human judgments, but it also introduces dependence on
the metric model and its training distribution. Under domain shift, a metric
trained or tuned on one distribution may no longer rank outputs reliably in a
new domain. Model bias and the use of a reference remain relevant limitations.

\subsection{Source-aware and semantic evaluation}

Reference-based metrics treat the reference as the principal target. In D2T,
the structured source should also be treated as an evaluation object. This is
especially important when a reference contains information that is not present
in the input or when a candidate expresses an input fact using a different
wording.

PARENT (Precision And Recall of Entailed N-grams from the Table) addresses this
problem for table-to-text generation~\cite{dhingra-etal-2019-handling}. It
compares the candidate with both the reference and the source table. Its
precision and recall calculations use entailment estimates to reward information
supported by the table and to reduce the effect of divergent reference text.
PARENT illustrates the value of task-specific metrics, although its original
formulation is designed for semi-structured tables and requires a suitable
source-to-text entailment procedure.

Dušek and Kasner~\cite{dusek-kasner-2020-evaluating} proposed an NLI-based
semantic-accuracy metric for D2T. The structured input is converted into simple
textual statements, after which natural-language inference is used in both
directions. Testing whether the input entails the output helps detect
unsupported information, while testing whether the output entails each input
fact helps detect omissions. The resulting labels can distinguish faithful
outputs, omissions, hallucinations, and combinations of these errors. This
approach directly targets the semantic requirements of D2T, but its accuracy
depends on the data-to-text conversion, the NLI model, and the assumptions about
whether all input facts should be realized.

\subsection{LLM-based judging}

An LLM can evaluate a generated text against the structured input under an
explicit set of criteria. It can act as a judge
whose behavior depends on its training, prompting, and decoding
configuration. Such judges can assess faithfulness, coverage, fluency, or
formatting without requiring lexical overlap with a reference, but their scores
remain model judgments and should be reported with their prompt and aggregation
procedure.

\subsection{Human evaluation}

Human evaluation remains necessary because automatic scores are proxies for
quality. Annotators can assess the generated text directly or judge it together
with the input representation, reference, or both. Typical criteria include
grammaticality, fluency, coherence, informativeness, coverage, repetition, and
factual consistency. For D2T, semantic criteria are often operationalized
through questions such as whether required facts are present and whether
unsupported facts have been added. Text quality can also be measured on a Likert
scale~\cite{gatt2018survey,novikova-etal-2017-need}.

Human evaluation is expensive and can be affected by annotator expertise,
instructions, sample selection, and the number of annotators per item. Studies
with multiple annotators should report inter-annotator agreement and describe
the assumptions behind the selected statistic. The survey by Artstein and
Poesio~\cite{artstein-poesio-2008-survey} discusses the use of agreement
coefficients in computational linguistics. Cohen's kappa measures
chance-corrected agreement between two annotators assigning nominal categories,
whereas Krippendorff's alpha is designed to accommodate multiple annotators,
missing annotations, and different measurement levels
\cite{cohen1960agreement,krippendorff2004reliability}. A single-rater study
cannot estimate agreement and should be presented as an exploratory assessment;
the manual pipeline assessment in Chapter~5 has this limitation. A reliable
evaluation should specify the criteria, scale, annotation procedure, number of
annotators, agreement measure where applicable, and aggregation method. Human
assessment is particularly important when automatic metrics disagree or when the
task permits many valid verbalizations
\cite{wang-etal-2023-faithful,xiang-etal-2022-asdot}.
